\documentclass[12pt,letterpaper]{article}

%% ─── Packages ──────────────────────────────────────────────────────────────
\usepackage[utf8]{inputenc}
\usepackage[T1]{fontenc}
\usepackage{lmodern}
\usepackage[margin=1.1in]{geometry}
\usepackage{amsmath,amssymb,amsthm}
\usepackage{mathtools}
\usepackage{microtype}
\usepackage{hyperref}
\usepackage{xcolor}
\usepackage{listings}
\usepackage{booktabs}
\usepackage{array}
\usepackage{longtable}
\usepackage{enumitem}
\usepackage{titlesec}
\usepackage{fancyhdr}
\usepackage{graphicx}
\usepackage{float}
\usepackage{mdframed}
\usepackage{tcolorbox}
\usepackage{datetime2}
\usepackage{setspace}
\usepackage{parskip}
\usepackage{etoolbox}

%% ─── Hyperref ───────────────────────────────────────────────────────────────
\hypersetup{
  colorlinks=true,
  linkcolor=black!70,
  urlcolor=teal!70!black,
  citecolor=black!60,
  pdftitle={PhaseMirror-HQ: Governance Protocol System — Prior Art Defensive Publication},
  pdfauthor={Ryan Van Gelder / Multiplicity Foundation},
  pdfsubject={Multiplicity Theory, PETC, ACE, Governance Protocols},
}

%% ─── Colors ─────────────────────────────────────────────────────────────────
\definecolor{primeblue}{RGB}{0,100,120}
\definecolor{codebackground}{RGB}{250,250,248}
\definecolor{codeborder}{RGB}{210,210,200}
\definecolor{boxaccent}{RGB}{1,105,111}
\definecolor{defcolor}{RGB}{40,80,100}
\definecolor{notebox}{RGB}{245,248,245}
\definecolor{warnbox}{RGB}{250,245,240}

%% ─── Section titles ─────────────────────────────────────────────────────────
\titleformat{\section}{\large\bfseries\color{primeblue}}{{\thesection}}{0.8em}{}[\vspace{2pt}\hrule\vspace{4pt}]
\titleformat{\subsection}{\normalsize\bfseries\color{black!80}}{{\thesubsection}}{0.7em}{}
\titleformat{\subsubsection}{\normalsize\itshape\color{black!70}}{{\thesubsubsection}}{0.6em}{}

%% ─── Headers/footers ────────────────────────────────────────────────────────
\pagestyle{fancy}
\fancyhf{}
\fancyhead[L]{\small\color{gray}\textit{PhaseMirror-HQ: Prior Art Defensive Publication}}
\fancyhead[R]{\small\color{gray}Citizen Gardens \textbullet{} 2026}
\fancyfoot[C]{\small\color{gray}\thepage}
\renewcommand{\headrulewidth}{0.3pt}
\renewcommand{\footrulewidth}{0pt}

%% ─── Listings ───────────────────────────────────────────────────────────────
\lstdefinestyle{pythonstyle}{
  language=Python,
  basicstyle=\ttfamily\small,
  backgroundcolor=\color{codebackground},
  frame=single,
  rulecolor=\color{codeborder},
  keywordstyle=\color{primeblue}\bfseries,
  commentstyle=\color{gray}\itshape,
  stringstyle=\color{teal!70!black},
  numberstyle=\tiny\color{gray},
  numbers=left,
  breaklines=true,
  showstringspaces=false,
  tabsize=4,
  xleftmargin=1.5em,
  framexleftmargin=1.5em,
  aboveskip=8pt,
  belowskip=8pt,
}

\lstdefinestyle{jsonlstyle}{
  language=,
  basicstyle=\ttfamily\footnotesize,
  backgroundcolor=\color{codebackground},
  frame=single,
  rulecolor=\color{codeborder},
  breaklines=true,
  showstringspaces=false,
  tabsize=2,
  xleftmargin=1em,
  aboveskip=6pt,
  belowskip=6pt,
}

\lstdefinestyle{circomstyle}{
  language=,
  basicstyle=\ttfamily\footnotesize,
  backgroundcolor=\color{codebackground},
  frame=single,
  rulecolor=\color{codeborder},
  keywordstyle=\color{primeblue}\bfseries,
  commentstyle=\color{gray}\itshape,
  breaklines=true,
  showstringspaces=false,
  tabsize=2,
  xleftmargin=1em,
  aboveskip=6pt,
  belowskip=6pt,
}

%% ─── Theorem environments ───────────────────────────────────────────────────
\theoremstyle{definition}
\newtheorem{definition}{Definition}[section]
\newtheorem{theorem}{Theorem}[section]
\newtheorem{proposition}{Proposition}[section]
\newtheorem{invariant}{Invariant}[section]
\theoremstyle{remark}
\newtheorem*{remark}{Remark}
\newtheorem*{notation}{Notation}

%% ─── Custom boxes ───────────────────────────────────────────────────────────
\tcbuselibrary{breakable,skins}
\newtcolorbox{priorartbox}{
  colback=notebox,
  colframe=boxaccent,
  fonttitle=\bfseries\small,
  title={Prior Art Disclosure Statement},
  breakable,
  arc=3pt,
  boxrule=0.5pt,
}
\newtcolorbox{claimbox}{
  colback=warnbox,
  colframe=primeblue!60,
  fonttitle=\bfseries\small,
  title={Technical Claim},
  breakable,
  arc=3pt,
  boxrule=0.5pt,
}
\newtcolorbox{noteboxtc}{
  colback=codebackground,
  colframe=codeborder,
  breakable,
  arc=2pt,
  boxrule=0.4pt,
  left=6pt, right=6pt, top=4pt, bottom=4pt,
}

%% ─── Spacing ────────────────────────────────────────────────────────────────
\setlength{\parskip}{6pt}
\setlength{\parindent}{0pt}
\onehalfspacing

%% ─────────────────────────────────────────────────────────────────────────────
\begin{document}

%% ══════════════════════════════════════════════════════════════════════════════
%%  TITLE PAGE
%% ══════════════════════════════════════════════════════════════════════════════
\begin{titlepage}
\centering

{\fontsize{11}{13}\selectfont\color{gray}\scshape
  Prior Art Defensive Publication \textbullet{} Technical Report}\\[0.6cm]

{\fontsize{26}{30}\selectfont\bfseries\color{primeblue}
  PhaseMirror-HQ\\[0.2cm]
  Governance Protocol System}\\[0.4cm]

{\fontsize{14}{17}\selectfont\color{black!75}
  Atomic Language Processing (ALP), PETC Exponent Vectors,\\[0.1cm]
  ACE Invariant Checking, and Multi-Layer Agent Architecture}\\[1.2cm]

\hrule height 0.5pt
\vspace{0.5cm}

\begin{minipage}{0.75\textwidth}
\centering
{\large\bfseries Citizen Gardens}\\[0.2cm]
{\normalsize The Foundation of Multiplicity}\\[0.2cm]
\end{minipage}

\vspace{0.5cm}
\hrule height 0.5pt
\vspace{1.0cm}

{\normalsize\color{gray}
  Publication Date: \DTMdate{2026-04-25}\\[0.2cm]
  Document Version: 1.0.0\\[0.2cm]
  Schema Version: 2.0.0 (index.json aligned)\\[0.2cm]
  Repository SHA Reference: \texttt{c348d1568af993f3ff26f10bd0ca470cb051bd06}
}

\vspace{1.2cm}

\begin{mdframed}[linecolor=boxaccent,linewidth=0.5pt,backgroundcolor=notebox,
  innerleftmargin=10pt,innerrightmargin=10pt,innertopmargin=8pt,innerbottommargin=8pt]
\small\textbf{Defensive Publication Notice.} This document constitutes a prior art
defensive publication for the purpose of placing all described methods, architectures,
mathematical frameworks, and software patterns into the public domain. The intent is to
prevent third-party patent claims on these techniques and to establish the Multiplicity
Foundation's priority date for all concepts herein. All technical content is released
under \textbf{Apache 2.0} (code) and \textbf{CC BY 4.0} (documentation) unless otherwise
noted. The Prime Materia License governs the core Multiplicity Theory IP layer.
\end{mdframed}

\vfill

{\small\color{gray}\textit{
  Multiplicity Theory constitutes a prime-indexed, recursively stable mathematical
  framework wherein multiplicity—the recurrence of elements understood through prime
  number decomposition—forms the basis for modeling nonlinear, emergent structures
  across mathematics, physics, computation, and cognition.}}

\end{titlepage}

%% ══════════════════════════════════════════════════════════════════════════════
%%  TABLE OF CONTENTS
%% ══════════════════════════════════════════════════════════════════════════════
\tableofcontents
\newpage

%% ══════════════════════════════════════════════════════════════════════════════
%%  EXECUTIVE SUMMARY
%% ══════════════════════════════════════════════════════════════════════════════
\section*{Executive Summary}
\addcontentsline{toc}{section}{Executive Summary}

This document is both a technical report and a prior art defensive publication covering
the complete five-phase architecture of PhaseMirror-HQ, a governance protocol system
developed under the Multiplicity Foundation. The system implements Atomic Language
Processing (ALP), a formal policy language grounded in Multiplicity Theory, wherein
every policy statement is decomposed into prime-indexed semantic dimensions and verified
through a stack of invariant-checking layers.

The five phases addressed are:
\begin{enumerate}[leftmargin=2em]
  \item \textbf{Phase 1 — ALP Grammar and Policy Pack Parsing} (Days 1–7): Lark-based
        grammar, PolicyTransformer, and JSON Schema validation.
  \item \textbf{Phase 2 — Predicates and PETC Binding} (Days 8–14): Token-to-predicate
        mapping, \texttt{predicates.py} module, and engine extension for predicate gating.
  \item \textbf{Phase 3 — Audit Layer} (Days 14–21): Structured JSON event emission,
        \texttt{audit.py}, tamper-evident fingerprinting, and circom ZK-SNARK binding.
  \item \textbf{Phase 4 — State Transitions and Governance Layer} (Days 21–30):
        \texttt{transitions.py} state machine, \texttt{index.json} versioning, exception
        workflow, and kill-switch in \texttt{manager.py}.
  \item \textbf{Phase 5 — Tooling and Agent Layer} (Days 30+): Bounded-scope agent
        architecture, \texttt{agents/} directory, \texttt{HandoffEnvelope} protocol, and
        ACE-constrained rank governance.
\end{enumerate}

The mathematical backbone is the PETC (Prime-Exponent Tensor Calculus) framework,
which maps every policy semantic dimension to an exponent vector \(\sigma \in
\mathbb{Z}^n\), and the ACE (Arithmetic Control Engine), which enforces spectral drift
bounds across update cycles. All code artifacts are hosted at \texttt{MultiplicityFoundation/PhaseMirror-HQ}.

\newpage

%% ══════════════════════════════════════════════════════════════════════════════
%%  1. BACKGROUND AND MOTIVATION
%% ══════════════════════════════════════════════════════════════════════════════
\section{Background and Motivation}

\subsection{Multiplicity Theory: Core Postulate}

Multiplicity Theory (MT) reinterprets mathematical objects not as static entities but
as recursively generated patterns of prime-labeled interactions. The foundational
postulate is:

\begin{definition}[Multiplicity Space]
A \emph{multiplicity space} \(\mathcal{M} = (S, \mathbf{p}, \phi)\) consists of a set
\(S\) of elements, a prime-index assignment \(\mathbf{p}: S \to \mathbb{P}^n\) (where
\(\mathbb{P}\) denotes the primes), and a recursive feedback operator
\(\phi: \mathcal{M} \times \mathcal{M} \to \mathcal{M}\) that preserves identity across
scales.
\end{definition}

Within MT, the cardinality of a relation is not simply a count but a prime-factored
multiplicity weight. For an association \(A \subseteq S \times S\), the multiplicity
at index \(k\) is:
\[
  \mu_k(A) = \prod_{(s_i, s_j) \in A} p_k^{e_{ij}} \tag{1}
\]
where \(p_k\) is the \(k\)-th prime and \(e_{ij}\) is the exponent encoding the
relational weight between \(s_i\) and \(s_j\).

\subsection{The Governance Motivation}

Policy languages in current practice (XACML, OPA/Rego, Cedar) treat policy evaluation
as a tree walk over a static document. They provide no formal guarantee that:
\begin{itemize}
  \item semantic dimensions are stable under recursive policy composition;
  \item audit trails are cryptographically linked to the evaluated state;
  \item state transitions in policy lifecycle are formally verified against invariants;
  \item agent-layer execution is bounded by the same mathematical constraints that
        govern the policy itself.
\end{itemize}

PhaseMirror-HQ addresses all four gaps through the Multiplicity Theory stack, with PETC
as the canonical state representation and ACE as the invariant-enforcement mechanism.

\subsection{Repository Context}

The implementation resides at \texttt{MultiplicityFoundation/PhaseMirror-HQ},
anchored to commit \texttt{c348d1568af993f3ff26f10bd0ca470cb051bd06}. The primary
library files relevant to this publication are:

\begin{table}[H]
\centering
\small
\begin{tabular}{@{}lll@{}}
\toprule
\textbf{File} & \textbf{Role} & \textbf{Phase} \\
\midrule
\texttt{multiplicity.py} & PETC exponent vector logic, \texttt{verify\_system} & 1, 2 \\
\texttt{association.py} & \texttt{PrimeAuditedAssociation}, multiplicity bounds & 2, 5 \\
\texttt{engine.py} & \texttt{EvolutionEngine}, entropy gating, evolve() & 2 \\
\texttt{manager.py} & Orchestration surface, link/audit/kill-switch & 3, 4 \\
\texttt{utils.py} & \texttt{get\_prime\_tag}, \texttt{multiplicity\_gate} & 2, 5 \\
\texttt{association\_verify.circom} & zk-SNARK for association integrity & 3 \\
\texttt{index.json} & Concept/article registry, versioning & 4 \\
\texttt{main.py} & Entry point, wiring pattern & 5 \\
\bottomrule
\end{tabular}
\caption{Core library files and their governance roles.}
\end{table}

\newpage

%% ══════════════════════════════════════════════════════════════════════════════
%%  2. MATHEMATICAL FRAMEWORK
%% ══════════════════════════════════════════════════════════════════════════════
\section{Mathematical Framework}

\subsection{PETC: Prime-Exponent Tensor Calculus}

\subsubsection{Exponent Vector State}

The canonical state of a policy pack \(\Pi\) is an integer-valued exponent vector:
\[
  \sigma(\Pi) = (\sigma_1, \sigma_2, \ldots, \sigma_n) \in \mathbb{Z}^n \tag{2}
\]
where each dimension \(\sigma_i\) corresponds to a semantic feature dimension \(q_i\)
(subject scope, action class, data class, constraint, exception authority). The prime
encoding of the state is:
\[
  \hat{\Pi} = \prod_{i=1}^{n} p_i^{\sigma_i} \tag{3}
\]
where \(p_i\) is the \(i\)-th prime. This encoding is injective by the Fundamental
Theorem of Arithmetic: distinct exponent vectors map to distinct rational numbers,
providing a canonical fingerprint for any policy state.

\subsubsection{Update Dynamics}

Let \(\sigma^{(t)}\) denote the PETC state at update cycle \(t\). An update
\(\Delta\sigma \in \mathbb{Z}^n\) is \emph{lawful} if and only if:
\[
  \|\sigma^{(t+1)} - \sigma^{(t)}\|_2 \leq \varepsilon \tag{4}
\]
for a drift bound \(\varepsilon > 0\) established at policy initialization. This is the
ACE spectral drift constraint, implemented as \texttt{ConstraintInvariantPredicate} in
\texttt{predicates.py}.

\subsubsection{Multiplicity Gate}

The \texttt{multiplicity\_gate} function in \texttt{utils.py} implements a prime-modular
projection:
\[
  \Gamma_{p_k}(\tau) = \tau \bmod p_k \tag{5}
\]
where \(\tau\) is a tensor state hash and \(p_k = \texttt{get\_prime\_tag}(k)\). This
gate modulates recursive link state per the PIRTM (Prime Interval Recursion Theory
Machine) specification, ensuring that state evolution is bounded within the prime
arithmetic of the \(k\)-th index.

\subsection{ACE: Arithmetic Control Engine}

\subsubsection{Invariant Checking Layer}

The ACE layer enforces three categories of invariants across every PETC update:

\begin{invariant}[Scope Membership]
For all subjects \(u \in \text{subject}(\Pi)\), the order prime \(p_u\) must be a
member of the declared scope set \(P_{\text{ord}}\):
\[
  p_u \in P_{\text{ord}} \quad \forall u \in \text{subject}(\Pi) \tag{6}
\]
\end{invariant}

\begin{invariant}[Feature Dimension Access]
For all actions \(a \in \text{action}(\Pi)\), the feature prime \(q_a\) must be in the
permitted action set \(Q_{\text{permit}}\):
\[
  q_a \in Q_{\text{permit}} = \{p : p \text{ indexes a permitted action}\} \tag{7}
\]
\end{invariant}

\begin{invariant}[Spectral Drift Bound]
The \(\ell^2\) norm of the PETC state must not exceed \(\varepsilon\) at any update
cycle (Equation 4). Equivalently, no single dimension may exceed the bound without
a corresponding decay in another dimension, preserving the total ``energy'' of the
policy state.
\end{invariant}

\subsubsection{Association Matrix}

The association matrix \(A \in \mathbb{Z}^{m \times m}\) (implemented in
\texttt{association.py}) encodes multiplicity bounds between source and target class
instances. Its entries \(A_{ij}\) count the number of active links from instance \(i\)
to instance \(j\). ACE enforces:
\[
  \text{lb}(A_{ij}) \leq A_{ij} \leq \text{ub}(A_{ij}) \tag{8}
\]
where \(\text{lb}, \text{ub}\) are lower and upper bounds derived from
\texttt{MultiplicityBound} (ONE, ZERO\_TO\_ONE, ZERO\_TO\_MANY, ONE\_TO\_MANY).

The ACE spectral drift condition extends to the association matrix update as follows.
Let \(A^{(t)}\) be the association matrix at cycle \(t\). A rank update is lawful if:
\[
  \|A^{(t+1)} - A^{(t)}\|_2 \leq \varepsilon_A \tag{9}
\]
where \(\varepsilon_A\) is the per-association drift bound and \(\|\cdot\|_2\) is the
operator (spectral) norm. This prevents oscillatory rank behavior and is the formal
basis for ACE-constrained agent rank updates in Phase 5.

\subsection{The \(\Lambda^p\)-Archivum}

Every audit event is tagged with a prime \(p_k = \texttt{get\_prime\_tag}(k)\), forming
a sequence of prime-indexed evidence records. The \(\Lambda^p\)-Archivum is the ordered
log of all such records:
\[
  \Lambda^p = \{(e_1, p_{k_1}), (e_2, p_{k_2}), \ldots, (e_N, p_{k_N})\} \tag{10}
\]
Each pair \((e_i, p_{k_i})\) is a structured audit event with a canonical SHA-256
fingerprint \(h_i = \text{SHA256}(\text{canonical}(e_i))\). The integrity of the log
is verified by re-computing all \(h_i\) and checking that each matches the stored value.

\subsection{Circom Association Integrity}

The \texttt{association\_verify.circom} circuit implements a zk-SNARK over the
following constraint system. Let \(s, t, m_s, m_t, p\) be the circuit inputs
(\texttt{source\_count}, \texttt{target\_count}, \texttt{source\_mult},
\texttt{target\_mult}, \texttt{prime\_tag}). Define Boolean validity signals:
\[
  v_s = \bigwedge_{b \in \{0,1,2,3\}} \text{mult\_valid}(s, m_s, b), \quad
  v_t = \bigwedge_{b \in \{0,1,2,3\}} \text{mult\_valid}(t, m_t, b) \tag{11}
\]
\[
  \text{is\_prime}(p) = [p \in \{2, 3, 5, 7\}] \quad \text{(mock circuit)} \tag{12}
\]
\[
  \text{is\_valid} = v_s \cdot v_t \cdot \text{is\_prime}(p) \tag{13}
\]
The output \texttt{is\_valid} $\in \{0,1\}$ is committed to the audit log via
\texttt{emit\_circom\_proof} in \texttt{audit.py}, with a SHA-256 witness hash of the
signal input vector serving as the \texttt{circom\_proof} field.

\newpage

%% ══════════════════════════════════════════════════════════════════════════════
%%  3. PHASE 1: ALP GRAMMAR AND POLICY PACK PARSING
%% ══════════════════════════════════════════════════════════════════════════════
\section{Phase 1: ALP Grammar and Policy Pack Parsing (Days 1–7)}

\subsection{Atomic Language Processing: Design Rationale}

ALP (Atomic Language Processing) treats every policy statement as a set of atomic
semantic units, each assigned to a prime-indexed dimension. Unlike XACML's XML-tree
evaluation or OPA's unification-based Rego, ALP's evaluation model is arithmetic:
each atom resolves to a dimension index \(i\) and an exponent value \(\sigma_i\),
producing a PETC state vector.

The grammar is defined over five token classes:

\begin{table}[H]
\centering
\small
\begin{tabular}{@{}llll@{}}
\toprule
\textbf{Token} & \textbf{Semantic Dimension} & \textbf{PETC Index} & \textbf{Example Value} \\
\midrule
\texttt{subject}    & Who the policy applies to & \(\sigma_1\) & \texttt{"patient"} \\
\texttt{action}     & What operation is permitted & \(\sigma_2\) & \texttt{"read"} \\
\texttt{data\_class} & Which data category & \(\sigma_3\) & \texttt{"PHI"} \\
\texttt{constraint} & Boundary condition & \(\sigma_4\) & \texttt{"within-jurisdiction"} \\
\texttt{exception}  & Override authority & \(\sigma_5\) & \texttt{\{allowed\_by: "chief-privacy-officer"\}} \\
\bottomrule
\end{tabular}
\caption{ALP grammar token classes and their PETC dimension bindings.}
\end{table}

\subsection{Lark Grammar Specification}

\begin{noteboxtc}
\begin{lstlisting}[style=pythonstyle, caption={ALP grammar excerpt (multiplicity.lark)}, label={lst:grammar}]
// multiplicity.lark — ALP policy grammar
// Each rule maps to one semantic dimension in sigma

start: policy_pack+

policy_pack: "policy" PACK_ID "{"
             scope_clause
             control_clause
             "}"

scope_clause: "scope" "{" data_class_list "}"
data_class_list: "data_classes" ":" "[" (ATOM ",")* ATOM "]"

control_clause: "control" "{"
                subject_rule
                action_rule
                constraint_rule?
                exception_rule?
                "}"

subject_rule: "subject" ":" "[" (ATOM ",")* ATOM "]"
action_rule:  "action"  ":" ATOM
constraint_rule: "constraint" ":" QUOTED_STRING
exception_rule: "exception" "{" "allowed_by" ":" ATOM
                              "condition" ":" QUOTED_STRING "}"

PACK_ID: /[a-z][a-z0-9\-]*/
ATOM:    /[a-zA-Z][a-zA-Z0-9_\-]*/
QUOTED_STRING: /"[^"]*"/
\end{lstlisting}
\end{noteboxtc}

\subsection{PolicyTransformer}

The \texttt{PolicyTransformer} class extends Lark's \texttt{Transformer} to convert
parse trees into validated \texttt{policy\_pack} dicts. Each method corresponds to one
grammar rule; the output is a Python dict conforming to \texttt{policy\_pack.schema.json}.

\begin{claimbox}
\textbf{Claim 1.} The PolicyTransformer implements a structure-preserving map
\(T: \mathcal{G} \to \mathcal{P}\) from grammar parse trees \(\mathcal{G}\) to
policy pack dicts \(\mathcal{P}\), such that every valid tree maps to a unique dict and
every dict admits a canonical PETC state vector \(\sigma \in \mathbb{Z}^5\).
\end{claimbox}

\subsection{JSON Schema Validation}

Policy packs are validated against \texttt{policy\_pack.schema.json} before entering
the predicate layer. The schema enforces:
\begin{itemize}
  \item \texttt{pack\_id}: non-empty string
  \item \texttt{scope.data\_classes}: array of non-empty strings
  \item \texttt{control.subject}: array of non-empty strings
  \item \texttt{control.action}: one of the permitted action set
  \item \texttt{control.constraint}: optional string
  \item \texttt{control.exception}: optional object with \texttt{allowed\_by} and
        \texttt{condition}
\end{itemize}

\newpage

%% ══════════════════════════════════════════════════════════════════════════════
%%  4. PHASE 2: PREDICATES AND PETC BINDING
%% ══════════════════════════════════════════════════════════════════════════════
\section{Phase 2: Predicates and PETC Binding (Days 8–14)}

\subsection{Predicate Architecture}

A \emph{predicate} \(P_i\) is a function:
\[
  P_i: \mathcal{P} \times \mathbb{Z}^n \to \{0,1\} \times \Sigma^* \times \mathbb{P}
  \tag{14}
\]
taking a policy pack \(\pi \in \mathcal{P}\) and PETC state \(\sigma \in \mathbb{Z}^n\),
and returning a tuple \((\text{passed}, \text{reason}, p_k)\) where \(\text{passed}
\in \{0,1\}\), \(\text{reason} \in \Sigma^*\) is a human-readable explanation, and
\(p_k\) is the prime tag from \texttt{utils.py}.

\subsection{Token-to-Predicate Mapping}

\begin{table}[H]
\centering
\small
\begin{tabular}{@{}lllp{4.5cm}@{}}
\toprule
\textbf{Grammar Token} & \textbf{Predicate Class} & \textbf{Invariant} & \textbf{PETC Binding} \\
\midrule
\texttt{subject}    & \texttt{SubjectScopePredicate}     & Subject \(\subseteq\) declared scope   & Order prime \(p_u\) membership in \(P_{\text{ord}}\) \\
\texttt{action}     & \texttt{ActionPermitPredicate}     & Action \(\in Q_{\text{permit}}\)       & Feature dimension \(q_i\) access \\
\texttt{data\_class} & \texttt{DataClassPredicate}       & \(\sigma_i \geq 0\), index mapped      & \(\sigma_i\) dimension bounds \\
\texttt{constraint} & \texttt{ConstraintInvariantPredicate} & \(\|\sigma\|_2 \leq \varepsilon\)   & ACE spectral drift (Eq.\ 4) \\
\texttt{exception}  & \texttt{ExceptionAuthorityPredicate} & Authority named; churn \(\leq K\)    & Bounded exception count per update \\
\bottomrule
\end{tabular}
\caption{Complete token-to-predicate mapping for the ALP grammar.}
\end{table}

\subsection{PredicateRegistry and Coverage Metric}

The \texttt{PredicateRegistry} maintains a dict of registered predicates keyed by
\texttt{predicate\_id}. The coverage metric is:
\[
  \text{coverage} = \frac{|\{P_i.\text{token\_type}\} \cap T|}{|T|} \tag{15}
\]
where \(T = \{\text{subject}, \text{action}, \text{data\_class}, \text{constraint},
\text{exception}\}\) is the full grammar token set. Coverage must equal 1.0 before
Phase 3 begins.

\subsection{Engine Extension: Predicate Gate}

The \texttt{EvolutionEngine.evolve()} method in \texttt{engine.py} is extended with
a predicate gate that fires before committing any PETC state update:

\begin{noteboxtc}
\begin{lstlisting}[style=pythonstyle, caption={Predicate gate in EvolutionEngine.evolve()}, label={lst:engine}]
def evolve(self, system_state, petc_state=None, policy_pack=None):
    # Existing entropy gate
    entropy = self.compute_entropy(system_state)
    if entropy > 1.0:
        raise ValueError("Entropy exceeds lawful threshold")

    # Phase 2: predicate gate
    if (self.predicate_registry is not None
            and petc_state is not None
            and policy_pack is not None):
        results = self.predicate_registry.run_all(policy_pack, petc_state)
        self.predicate_log.append(results)
        failures = [r for r in results if not r.passed]
        if failures:
            reasons = "; ".join(r.reason for r in failures)
            raise ValueError(f"Predicate gate failed: {reasons}")

    modulated = multiplicity_gate(self.prime_index, system_state)
    self.history.append(modulated)
    self.prime_index += 1
    return modulated
\end{lstlisting}
\end{noteboxtc}

\newpage

%% ══════════════════════════════════════════════════════════════════════════════
%%  5. PHASE 3: AUDIT LAYER
%% ══════════════════════════════════════════════════════════════════════════════
\section{Phase 3: Audit Layer (Days 14–21)}

\subsection{Design Principles}

Three principles govern the audit layer design:
\begin{enumerate}
  \item \textbf{Every control event is a record.} No state transition in
        \texttt{manager.py}, \texttt{engine.py}, or \texttt{predicates.py} may execute
        without emitting an \texttt{AuditEvent}.
  \item \textbf{Every record has a fingerprint.} SHA-256 of the canonical JSON
        serialization provides tamper evidence without requiring a full Merkle tree.
  \item \textbf{Every fingerprint is a commitment.} Fingerprints may be committed to
        the \texttt{circom\_proof} field for association integrity events, creating a
        chain of cryptographic evidence linked to \texttt{association\_verify.circom}.
\end{enumerate}

\subsection{Audit Event Schema}

\begin{noteboxtc}
\begin{lstlisting}[style=jsonlstyle, caption={AuditEvent JSON schema (schema version 1.0.0)}, label={lst:auditevent}]
{
  "event_id":       "uuid4",
  "timestamp":      "ISO-8601 (UTC)",
  "event_type":     "define_association | link | predicate_check |
                     rank_swap | petc_update | audit_pass |
                     audit_fail | exception_granted | kill_switch",
  "policy_id":      "string | null",
  "control": {
    "subject":      ["string"],
    "action":       "string",
    "data_class":   ["string"]
  },
  "result": {
    "passed":       "boolean",
    "reason":       "string",
    "prime_tag":    "integer (prime)"
  },
  "owner":          "string",
  "schema_version": "1.0.0",
  "circom_proof":   "hex-string | null"
}
\end{lstlisting}
\end{noteboxtc}

\subsection{AuditLog and Coverage Metric}

The \texttt{AuditLog} class maintains an append-only in-memory buffer and optionally
writes to a JSONL file. The audit-trace completeness metric is:
\[
  \text{completeness} = \frac{N_{\text{logged}}}{N_{\text{expected}}} \tag{16}
\]
where \(N_{\text{expected}}\) is the count of control event entry-points in
\texttt{manager.py} (incremented on each call to any public method) and
\(N_{\text{logged}}\) is the count of successfully emitted \texttt{AuditEvent}
instances. The target is \(\text{completeness} = 1.0\).

\subsection{Circom Binding}

The \texttt{emit\_circom\_proof} function mirrors the signal names of
\texttt{association\_verify.circom} exactly:

\begin{noteboxtc}
\begin{lstlisting}[style=circomstyle, caption={association\_verify.circom — signal interface}, label={lst:circom}]
template AssociationVerify() {
    signal input source_count;  // Number of source instances
    signal input target_count;  // Number of target instances
    signal input source_mult;   // Encoded multiplicity (0..3)
    signal input target_mult;
    signal input prime_tag;     // Prime tag for auditability
    signal output is_valid;     // 0 or 1

    // is_valid = source_valid * target_valid * is_prime(prime_tag)
}
\end{lstlisting}
\end{noteboxtc}

The witness hash \(h_w = \text{SHA256}(\text{canonical}(\{s_c, t_c, m_s, m_t, p, v\}))\)
is stored in \texttt{circom\_proof}. A full snarkjs proof generation path can replace
the hash in Phase 5+ without changing the audit schema.

\newpage

%% ══════════════════════════════════════════════════════════════════════════════
%%  6. PHASE 4: STATE TRANSITIONS AND GOVERNANCE LAYER
%% ══════════════════════════════════════════════════════════════════════════════
\section{Phase 4: State Transitions and Governance Layer (Days 21–30)}

\subsection{Gating Condition}

Phase 4 is gated on Phase 2 invariant stability: \(\text{coverage}(\mathtt{PredicateRegistry}) = 1.0\)
and invariant pass rate \(\geq 0.99\) across the test suite. This is not optional;
state transitions encode behavior over time and incorrect predicates produce
incorrect transitions that are hard to roll back.

\subsection{Policy Pack Lifecycle: State Machine}

A policy pack transitions through five lifecycle states:
\[
  \texttt{draft} \to \texttt{active} \to \texttt{suspended} \to \texttt{deprecated} \to \texttt{archived}
  \tag{17}
\]

The legal transition set \(\mathcal{T}\) is:
\[
  \mathcal{T} = \{(\texttt{draft}, \texttt{active}),\ (\texttt{active}, \texttt{suspended}),\
  (\texttt{active}, \texttt{deprecated}),\ (\texttt{suspended}, \texttt{active}),
\]
\[
  (\texttt{suspended}, \texttt{deprecated}),\ (\texttt{deprecated}, \texttt{archived})\}
  \tag{18}
\]

Every transition attempt \((s_{\text{from}}, s_{\text{to}})\) not in \(\mathcal{T}\) is
rejected with an \texttt{audit\_fail} event. Every successful transition emits
\texttt{audit\_pass}. The transition validity rate is:
\[
  \text{TVR} = \frac{|\{r \in \mathcal{H} : r.\text{allowed}\}|}{|\mathcal{H}|} \tag{19}
\]
where \(\mathcal{H}\) is the transition history of a \texttt{PolicyStateMachine} instance.

\subsection{Index.json Versioning}

The \texttt{index.json} file is extended with three governance fields at both root and
article level:

\begin{table}[H]
\centering
\small
\begin{tabular}{@{}llp{5cm}@{}}
\toprule
\textbf{Field} & \textbf{Type} & \textbf{Semantics} \\
\midrule
\texttt{schema\_version} & \texttt{"2.0.0"} & Version of the index.json schema itself \\
\texttt{deprecated\_at}  & ISO-8601 or \texttt{null} & When the article was deprecated; null if active \\
\texttt{successor\_id}   & slug string or \texttt{null} & The slug of the replacing article; null if none \\
\bottomrule
\end{tabular}
\caption{New governance fields added to index.json in Phase 4.}
\end{table}

The invariant is: \texttt{deprecated\_at} $\neq$ \texttt{null} $\Rightarrow$
\texttt{successor\_id} $\neq$ \texttt{null}. This ensures deprecation is always
paired with a migration path.

\subsection{Exception Workflow and Kill-Switch}

\subsubsection{Grant and Revoke}

The \texttt{grant\_exception()} method registers a bounded exception in
\texttt{manager.py}. The mean time to revoke (MTTR) is tracked as:
\[
  \text{MTTR} = \frac{1}{|\mathcal{R}|} \sum_{r \in \mathcal{R}} (t_r^{\text{revoke}} - t_r^{\text{grant}}) \tag{20}
\]
where \(\mathcal{R}\) is the set of revoked exceptions and times are in seconds.

\subsubsection{Kill-Switch}

The \texttt{kill\_switch()} method:
\begin{enumerate}
  \item Sets \texttt{\_kill\_switch\_active = True}, blocking all future exception grants.
  \item Iterates all \texttt{ACTIVE} policy packs in the \texttt{TransitionRegistry} and
        transitions each to \texttt{SUSPENDED} with a kill-switch reason string.
  \item Emits a single \texttt{kill\_switch} event to the audit log.
\end{enumerate}

The kill-switch-to-suspend latency must be \(< 1\) second in test, verified by
timing the \texttt{kill\_switch()} call across a registry with \(n = 1000\) active
packs.

\newpage

%% ══════════════════════════════════════════════════════════════════════════════
%%  7. PHASE 5: TOOLING AND AGENT LAYER
%% ══════════════════════════════════════════════════════════════════════════════
\section{Phase 5: Tooling and Agent Layer (Days 30+)}

\subsection{Agent Architecture Contract}

Each agent is a bounded execution unit satisfying three hard constraints:
\begin{enumerate}
  \item \textbf{Single read source}: reads from exactly one layer (L1, L2, or L3).
  \item \textbf{Single write artifact}: writes to exactly one typed output.
  \item \textbf{Audit on every action}: every \texttt{execute()} call emits at least
        one \texttt{AuditEvent} regardless of success or failure.
\end{enumerate}

These constraints are enforced by \texttt{BaseAgent.run()}, which wraps every
\texttt{execute()} call in a try/except that emits \texttt{audit\_fail} on any
unhandled exception.

\subsection{Agent Pipeline}

\[
  \underbrace{\texttt{ParserAgent}}_{\text{L1}}
  \xrightarrow{\texttt{HandoffEnvelope}(\texttt{ParsedPack})}
  \underbrace{\texttt{PredicateAgent}}_{\text{L2}}
  \xrightarrow{\texttt{HandoffEnvelope}(\texttt{List[PredicateResult]})}
  \underbrace{\texttt{EvidenceAgent}}_{\text{L3}}
  \xrightarrow{\texttt{HandoffEnvelope}(\texttt{AuditBundle})}
  \texttt{audit.py}
  \tag{21}
\]

The \texttt{HandoffEnvelope} is the typed communication primitive between agents. It
carries a \texttt{passed} boolean; a failed envelope short-circuits the pipeline and
increments the receiving agent's \texttt{handoff\_failure\_rate}.

\subsection{ACE-Constrained Agent Ranking}

Agent rank across sessions is governed by \texttt{PrimeAuditedAssociation}
(\texttt{association.py}). Each agent is the source class in an association with its
task class. The ACE rank update rule (Eq.\ 9) applies: no agent may jump more than
\(K\) rank positions per update cycle. This prevents the oscillatory priority
inversion that afflicts naive round-robin agent schedulers.

\begin{definition}[Bounded Agent Rank]
Let \(r_a^{(t)} \in \mathbb{Z}^+\) be the rank of agent \(a\) at session \(t\). The
rank update is \emph{ACE-bounded} if:
\[
  |r_a^{(t+1)} - r_a^{(t)}| \leq K \quad \forall a, t \tag{22}
\]
where \(K\) is the bounded-churn constant, matching the \texttt{ExceptionAuthorityPredicate}
exception count limit established in Phase 2.
\end{definition}

\subsection{Agent Metrics}

\begin{table}[H]
\centering
\small
\begin{tabular}{@{}lll@{}}
\toprule
\textbf{Metric} & \textbf{Formula} & \textbf{Target} \\
\midrule
Agent success rate & \(\frac{N_{\text{success}}}{N_{\text{total}}}\) & \(\geq 0.95\) per agent \\
Handoff failure rate & \(\frac{N_{\text{handoff\_fail}}}{N_{\text{total}}}\) & \(\leq 0.02\) \\
Pipeline completion rate & Envelopes reaching EvidenceAgent / all started & \(\geq 0.90\) \\
Audit-trace completeness & \texttt{AuditLog.coverage()} across all agents & \(= 1.0\) \\
Agent rank churn & \(\max_{a,t} |r_a^{(t+1)} - r_a^{(t)}|\) & \(\leq K\) \\
Evidence bundle integrity & Fingerprint recompute match rate & \(= 1.0\) \\
\bottomrule
\end{tabular}
\caption{Phase 5 agent metrics, formulas, and targets.}
\end{table}

\newpage

%% ══════════════════════════════════════════════════════════════════════════════
%%  8. CROSS-PHASE METRICS SUMMARY
%% ══════════════════════════════════════════════════════════════════════════════
\section{Cross-Phase Metrics Summary}

\begin{longtable}{@{}lllll@{}}
\toprule
\textbf{Metric} & \textbf{Phase} & \textbf{Owner} & \textbf{Target} & \textbf{Source} \\
\midrule
\endfirsthead
\toprule
\textbf{Metric} & \textbf{Phase} & \textbf{Owner} & \textbf{Target} & \textbf{Source} \\
\midrule
\endhead
Predicate completeness        & 2 & Formalism  & 100\% (5/5)       & \texttt{PredicateRegistry.coverage()} \\
Invariant pass rate           & 2 & Formalism  & \(\geq 99\%\)     & Test suite \\
Audit-trace completeness      & 3 & Security   & 100\%             & \texttt{AuditLog.coverage()} \\
Fingerprint integrity         & 3 & Security   & 100\%             & SHA-256 recompute \\
Circom proof coverage         & 3 & Security   & 100\% of links    & \texttt{query("audit\_pass")} \\
Transition validity rate      & 4 & Platform   & \(\geq 99\%\)     & \texttt{TransitionRegistry.overall\_validity\_rate()} \\
MTTR exception                & 4 & Governance & Baseline est.     & \texttt{MultiplicityManager.mean\_time\_to\_revoke()} \\
Kill-switch latency           & 4 & Platform   & \(< 1\)s          & Timing test (\(n=1000\)) \\
Index deprecation coverage    & 4 & Governance & 100\%             & \texttt{index.json} audit \\
Agent success rate            & 5 & Engineering& \(\geq 95\%\)     & \texttt{BaseAgent.success\_rate()} \\
Handoff failure rate          & 5 & Engineering& \(\leq 2\%\)      & \texttt{BaseAgent.handoff\_failure\_rate()} \\
Agent rank churn              & 5 & Engineering& \(\leq K\)        & ACE rank delta \\
\bottomrule
\caption{Cross-phase metrics, owners, targets, and measurement sources.}
\end{longtable}

\newpage

%% ══════════════════════════════════════════════════════════════════════════════
%%  9. REPRESENTATIVE CODE ARTIFACTS
%% ══════════════════════════════════════════════════════════════════════════════
\section{Representative Code Artifacts}

\subsection{PredicateRegistry: Core Interface}

\begin{noteboxtc}
\begin{lstlisting}[style=pythonstyle, caption={PredicateRegistry — registration and coverage metric}, label={lst:registry}]
class PredicateRegistry:
    """Registry of all active predicates, keyed by predicate_id."""
    def __init__(self):
        self._predicates: Dict[str, BasePredicate] = {}

    def register(self, predicate: BasePredicate):
        self._predicates[predicate.predicate_id] = predicate

    def run_all(self, policy_pack: dict,
                petc_state: np.ndarray) -> List[PredicateResult]:
        return [p.check(policy_pack, petc_state)
                for p in self._predicates.values()]

    def coverage(self) -> float:
        expected = {"subject","action","data_class","constraint","exception"}
        covered  = {p.token_type for p in self._predicates.values()}
        return len(covered & expected) / len(expected)  # Eq. 15
\end{lstlisting}
\end{noteboxtc}

\subsection{PolicyStateMachine: Transition Guard}

\begin{noteboxtc}
\begin{lstlisting}[style=pythonstyle, caption={PolicyStateMachine.transition() with audit emission}, label={lst:statemachine}]
def transition(self, to_state: PolicyState,
               reason: str = "") -> TransitionResult:
    pair = (self._state, to_state)
    if pair not in LEGAL_TRANSITIONS:          # Eq. 18
        result = TransitionResult(allowed=False, from_state=self._state,
                                  to_state=to_state,
                                  reason=f"Illegal: {self._state}->{to_state}")
        self._emit(result)
        return result
    if self.gate and not self.gate(self.pack_id, to_state):
        result = TransitionResult(allowed=False, from_state=self._state,
                                  to_state=to_state,
                                  reason=f"Gate rejected {to_state}")
        self._emit(result)
        return result
    prev = self._state
    self._state = to_state
    result = TransitionResult(allowed=True, from_state=prev,
                              to_state=to_state, reason=reason)
    self._history.append(result)
    self._emit(result)  # Audit every transition — Eq. 16
    return result
\end{lstlisting}
\end{noteboxtc}

\subsection{BaseAgent: Bounded Execution}

\begin{noteboxtc}
\begin{lstlisting}[style=pythonstyle, caption={BaseAgent.run() — audit-on-every-action wrapper}, label={lst:baseagent}]
def run(self, envelope: HandoffEnvelope) -> HandoffEnvelope:
    self._total_count += 1
    # Propagate upstream failures without executing
    if not envelope.passed:
        self._handoff_failures += 1
        self._emit("audit_fail", envelope.task_id,
                   f"Received failed envelope from {envelope.source_agent}",
                   False)
        return HandoffEnvelope(source_agent=self.agent_id,
                               target_agent="none", payload=None,
                               task_id=envelope.task_id, passed=False,
                               failure_reason="Upstream failure")
    try:
        result = self.execute(envelope.payload, envelope.task_id)
        if result.success:
            self._success_count += 1
        self._emit("audit_pass" if result.success else "audit_fail",
                   envelope.task_id, result.reason, result.success)
        return HandoffEnvelope(source_agent=self.agent_id, target_agent="",
                               payload=result.output,
                               task_id=envelope.task_id,
                               passed=result.success,
                               failure_reason=None if result.success
                                              else result.reason)
    except Exception as e:
        self._emit("audit_fail", envelope.task_id, str(e), False)
        return HandoffEnvelope(source_agent=self.agent_id, target_agent="none",
                               payload=None, task_id=envelope.task_id,
                               passed=False, failure_reason=str(e))
\end{lstlisting}
\end{noteboxtc}

\subsection{AuditLog: Fingerprint and Coverage}

\begin{noteboxtc}
\begin{lstlisting}[style=pythonstyle, caption={AuditLog.emit() and AuditEvent.fingerprint()}, label={lst:auditlog}]
def emit(self, event: AuditEvent) -> str:
    """Append event, write to JSONL if path set. Returns fingerprint."""
    self._expected_count += 1
    self._events.append(event)
    self._logged_count += 1
    if self._log_path:
        with open(self._log_path, "a") as f:
            f.write(event.to_json() + "\n")
    return event.fingerprint()  # SHA-256 commitment

def fingerprint(self) -> str:
    """SHA-256 of canonical event for tamper detection."""
    canonical = json.dumps(self.to_dict(), sort_keys=True, default=str)
    return hashlib.sha256(canonical.encode()).hexdigest()
\end{lstlisting}
\end{noteboxtc}

\newpage

%% ══════════════════════════════════════════════════════════════════════════════
%%  10. PRIOR ART CLAIMS
%% ══════════════════════════════════════════════════════════════════════════════
\section{Prior Art Claims}

\begin{priorartbox}
All claims below are asserted as prior art as of the publication date of this document
(2026-04-25). This publication is intended to prevent the patenting of these techniques
by any party. Implementation code is available at:
\texttt{https://github.com/MultiplicityFoundation/PhaseMirror-HQ}.
\end{priorartbox}

\begin{claimbox}
\textbf{Claim A — Prime-Indexed Policy State Vectors (PETC).}\\
A method of representing policy semantics as an integer-valued exponent vector
\(\sigma \in \mathbb{Z}^n\) indexed by prime numbers, where each component corresponds
to a semantic dimension of a policy grammar token, and the canonical encoding is the
product \(\prod_{i=1}^n p_i^{\sigma_i}\) uniquely identifying the policy state by
the Fundamental Theorem of Arithmetic.
\end{claimbox}

\begin{claimbox}
\textbf{Claim B — ACE Spectral Drift Invariant.}\\
A method of enforcing policy stability through an \(\ell^2\)-norm drift constraint
\(\|\sigma^{(t+1)} - \sigma^{(t)}\|_2 \leq \varepsilon\) on consecutive PETC state
vectors, implemented as a predicate gate in a governance engine that rejects any policy
update violating this bound.
\end{claimbox}

\begin{claimbox}
\textbf{Claim C — Token-to-Predicate Binding Architecture.}\\
A governance system in which each grammar token class of a formal policy language is
bound to a corresponding predicate function that takes a policy pack and a PETC state
vector and returns a typed result tuple including a pass/fail Boolean, a human-readable
reason string, and a prime tag for audit tracing.
\end{claimbox}

\begin{claimbox}
\textbf{Claim D — Append-Only Prime-Tagged Audit Log with SHA-256 Commitments.}\\
A structured audit logging system wherein every control event in a governance engine
emits a JSON record with a \texttt{prime\_tag} field drawn from the prime-index
sequence, a \texttt{fingerprint} field computed as SHA-256 of the canonical
serialization of the event, and a \texttt{circom\_proof} field optionally holding a
cryptographic witness hash from a zk-SNARK circuit verifying association integrity.
\end{claimbox}

\begin{claimbox}
\textbf{Claim E — Policy Pack Lifecycle State Machine with Predicate Gate.}\\
A state machine governing the lifecycle of formal policy packs through the states
\{draft, active, suspended, deprecated, archived\} with a legal transition set
\(\mathcal{T}\) (Eq.\ 18), wherein every transition attempt is guarded by an external
predicate gate callable and every transition outcome emits a structured audit event.
\end{claimbox}

\begin{claimbox}
\textbf{Claim F — Kill-Switch with Guaranteed Audit Emission.}\\
A governance method that, upon invocation, atomically: (1) sets a kill-switch flag
blocking all future exception grants and state transitions, (2) transitions all
ACTIVE policy packs to SUSPENDED, and (3) emits a single \texttt{kill\_switch} audit
event with timestamp, owner, and reason, ensuring no state change occurs without
a corresponding audit record even under emergency conditions.
\end{claimbox}

\begin{claimbox}
\textbf{Claim G — Bounded-Scope Agent Architecture with HandoffEnvelope.}\\
A multi-agent execution architecture wherein each agent has a declared read-source layer
and write-artifact type, communicates exclusively through typed \texttt{HandoffEnvelope}
objects carrying a \texttt{passed} Boolean for pipeline short-circuiting, and emits an
audit event on every action including upstream failure propagation.
\end{claimbox}

\begin{claimbox}
\textbf{Claim H — ACE-Constrained Agent Rank Governance.}\\
A method of governing agent priority across sessions using the same
\texttt{PrimeAuditedAssociation} multiplicity-bound mechanism that governs policy
association integrity, wherein agent rank updates satisfy a bounded-churn constraint
\(|r_a^{(t+1)} - r_a^{(t)}| \leq K\) (Eq.\ 22) preventing oscillatory priority
inversion.
\end{claimbox}

\newpage

%% ══════════════════════════════════════════════════════════════════════════════
%%  11. FILE MANIFEST AND DEPENDENCY GRAPH
%% ══════════════════════════════════════════════════════════════════════════════
\section{File Manifest and Dependency Graph}

\subsection{Existing Files (at reference commit)}

\begin{table}[H]
\centering
\small
\begin{tabular}{@{}lll@{}}
\toprule
\textbf{File} & \textbf{SHA (reference commit)} & \textbf{Role} \\
\midrule
\texttt{multiplicity.py}             & \texttt{c348d1568a...} & PETC logic, verify\_system \\
\texttt{association.py}              & \texttt{c348d1568a...} & PrimeAuditedAssociation, bounds \\
\texttt{engine.py}                   & \texttt{c348d1568a...} & EvolutionEngine, evolve() \\
\texttt{manager.py}                  & SHA: \texttt{7c91dc53d3...} & Orchestration, link, audit \\
\texttt{utils.py}                    & SHA: \texttt{6d07d57f17...} & get\_prime\_tag, mult\_gate \\
\texttt{main.py}                     & SHA: \texttt{de86a3ff81...} & Entry point \\
\texttt{association\_verify.circom}  & SHA: \texttt{267d62122e...} & zk-SNARK circuit \\
\texttt{index.json}                  & SHA: \texttt{02a1649b0e...} & Concept registry \\
\bottomrule
\end{tabular}
\caption{Existing library files with retrieved SHAs.}
\end{table}

\subsection{New Files Introduced (Phases 2–5)}

\begin{table}[H]
\centering
\small
\begin{tabular}{@{}llll@{}}
\toprule
\textbf{File} & \textbf{Phase} & \textbf{Owner} & \textbf{Key Export} \\
\midrule
\texttt{predicates.py}               & 2 & Formalism  & \texttt{PredicateRegistry}, 5 predicate classes \\
\texttt{audit.py}                    & 3 & Security   & \texttt{AuditLog}, \texttt{AuditEvent}, \texttt{emit\_circom\_proof} \\
\texttt{transitions.py}              & 4 & Platform   & \texttt{PolicyStateMachine}, \texttt{TransitionRegistry} \\
\texttt{agents/\_\_init\_\_.py}      & 5 & Engineering & Agent re-exports \\
\texttt{agents/base\_agent.py}       & 5 & Engineering & \texttt{BaseAgent}, \texttt{HandoffEnvelope} \\
\texttt{agents/parser\_agent.py}     & 5 & Engineering & \texttt{ParserAgent} (L1) \\
\texttt{agents/predicate\_agent.py}  & 5 & Engineering & \texttt{PredicateAgent} (L2) \\
\texttt{agents/evidence\_agent.py}   & 5 & Engineering & \texttt{EvidenceAgent}, \texttt{AuditBundle} (L3) \\
\texttt{agents/orchestrator.py}      & 5 & Engineering & \texttt{AgentOrchestrator}, \texttt{metrics()} \\
\bottomrule
\end{tabular}
\caption{New files introduced across Phases 2–5.}
\end{table}

\subsection{Dependency Graph}

\begin{noteboxtc}
\begin{lstlisting}[style=jsonlstyle, caption={Module dependency graph (simplified)}, label={lst:deps}]
main.py
 └─ manager.py
     ├─ association.py  ←─── multiplicity.py
     ├─ engine.py       ←─── utils.py
     │                  ←─── predicates.py  ←─── utils.py
     ├─ audit.py        ←─── association_verify.circom (proof commitment)
     └─ transitions.py  ←─── audit.py

agents/
 └─ orchestrator.py
     ├─ parser_agent.py    ←─── (lark grammar, json schema)
     ├─ predicate_agent.py ←─── predicates.py, audit.py
     └─ evidence_agent.py  ←─── audit.py
         └─ AuditBundle    ←─── audit.py (JSONL sink)
\end{lstlisting}
\end{noteboxtc}

\newpage

% ---- Title block ----
\appendix
This appendix collects the formal definitions, explicit proofs, and operator norm bounds underlying the PhaseMirror-HQ governance system.
The material proceeds from foundational number-theoretic constructions through
the Prime Exponent Transform over Constraint-space (\(\PETC\)), the
Arithmetic Control Engine (\(\ACE\)) spectral drift operator, and the
Multiplicity-Knot-Theoretic (\(\MKT\)) composition functor.
Each section corresponds to a module in the codebase:
\texttt{mkt\_sigma.py}, \texttt{omega\_operator.py},
\texttt{formal\_verification.py}, \texttt{galois\_certification.py},
\texttt{debug\_spectral.py}, and \texttt{lambda\_m\_protocols.py}.
All operator norm bounds are stated with explicit constants; proofs are
constructive where possible to support algorithmic extraction.

% ============================================================
\section{Notation and Foundational Conventions}
\label{sec:notation}
% ============================================================

Let \(\primes = \{p_1, p_2, p_3, \ldots\} = \{2, 3, 5, \ldots\}\) denote the
set of rational primes in ascending order.  For \(n \in \N\) write
\(\omega(n)\) for the number of distinct prime divisors and \(\Omega(n)\) for
the number of prime divisors counted with multiplicity.

\begin{definition}[Prime Index Map]
  \label{def:prime-index}
  The \emph{prime index map} \(\pi : \N \setminus \{0\} \to \N_0^{(\infty)}\)
  assigns to each positive integer \(n\) its exponent vector
  \[
    \pi(n) \;=\; \bigl(v_{p_1}(n),\, v_{p_2}(n),\, v_{p_3}(n),\,\ldots\bigr),
  \]
  where \(v_p(n)\) is the \(p\)-adic valuation of \(n\).  The vector
  \(\pi(n)\) has finite support (finitely many nonzero entries).
\end{definition}

\begin{definition}[Multiplicity Space]
  \label{def:mspace}
  A \emph{multiplicity space of order} \(k\) is the free \(\Z_{\ge 0}\)-module
  \[
    \Mspace{k} \;=\; \Z_{\ge 0}^k,
  \]
  equipped with component-wise partial order \(\le\) and the product monoid
  operation \(\oplus\) (coordinate-wise addition).  Elements
  \(\bm{e} \in \Mspace{k}\) are called \emph{exponent vectors}.
\end{definition}

Throughout, all vectors are column vectors; \(\bm{e}^\top\) denotes the
transpose.  For a matrix \(A\), \(\norm{A}_2\) is the spectral norm
(largest singular value), \(\norm{A}_F\) is the Frobenius norm, and
\(\sprad(A)\) is the spectral radius.

% ============================================================
\section{The PETC Transform}
\label{sec:petc}
% ============================================================

\subsection{Definition}
\label{ssec:petc-def}

Let \(\mathcal{C}\) be a finite constraint alphabet of cardinality \(m\).
A \emph{policy document} \(D\) is a sequence
\((c_1, c_2, \ldots, c_\ell)\) with each \(c_i \in \mathcal{C}\).

\begin{definition}[PETC Encoding]
  \label{def:petc}
  Fix an injective labelling \(\lambda : \mathcal{C} \to \primes\) mapping
  each constraint token to a distinct prime.  The
  \emph{Prime Exponent Transform over Constraint-space} of \(D\) is
  \[
    \PETC(D) \;=\; \pi\!\left(\prod_{i=1}^{\ell} \lambda(c_i)\right)
    \;=\; \sum_{c \in \mathcal{C}} \#\{i : c_i = c\}\cdot\bm{e}_{\lambda(c)},
  \]
  where \(\bm{e}_p\) is the standard basis vector for prime \(p\) and the sum
  is in \(\Mspace{|\mathcal{C}|}\).
\end{definition}

\noindent
Because \(\lambda\) is injective and primes are multiplicatively independent,
the Fundamental Theorem of Arithmetic immediately yields:

\begin{proposition}[PETC Injectivity]
  \label{prop:petc-inj}
  The map \(\PETC : \mathcal{C}^* \to \Mspace{|\mathcal{C}|}\) is injective
  on the set of \emph{multisets} over \(\mathcal{C}\).
\end{proposition}

\begin{proof}
  Two sequences with the same multiset of tokens produce identical products
  under \(\lambda\).  Conversely, if
  \(\prod_i \lambda(c_i) = \prod_j \lambda(c_j')\), then by unique
  factorisation and injectivity of \(\lambda\), the prime factorisation of
  both products is identical, so the multisets of tokens agree.
\end{proof}

\subsection{Metric Structure on PETC Space}
\label{ssec:petc-metric}

\begin{definition}[PETC Distance]
  \label{def:petc-dist}
  For two exponent vectors \(\bm{u}, \bm{v} \in \Mspace{k}\), define
  \[
    d_{\PETC}(\bm{u}, \bm{v})
    \;=\; \norm{\bm{u} - \bm{v}}_1
    \;=\; \sum_{i=1}^{k} \abs{u_i - v_i}.
  \]
\end{definition}

\begin{proposition}[Triangle Inequality]
  \(d_{\PETC}\) is a metric on \(\Z^k\), hence on \(\Mspace{k} \subset \Z^k\).
\end{proposition}

\begin{proof}
  Follows immediately from the \(\ell^1\)-norm triangle inequality on \(\Z^k\).
\end{proof}

\subsection{PETC Constraint Composition}
\label{ssec:petc-compose}

\begin{definition}[Policy Join]
  \label{def:policy-join}
  For documents \(D_1, D_2\), define
  \[
    D_1 \sqcup D_2 \;:=\; \PETC^{-1}\!\bigl(\PETC(D_1) \oplus \PETC(D_2)\bigr),
  \]
  the \emph{multiset union} policy.
\end{definition}

\begin{theorem}[Associativity and Commutativity of Join]
  \label{thm:join-assoc}
  The operation \(\sqcup\) is associative and commutative on policy documents.
\end{theorem}

\begin{proof}
  Both properties are inherited from coordinate-wise addition in
  \(\Mspace{k}\), which is a commutative monoid.
\end{proof}

\begin{theorem}[PETC Subsumption Order]
  \label{thm:petc-subsumption}
  Define \(D_1 \preceq D_2\) iff \(\PETC(D_1) \le \PETC(D_2)\) component-wise
  (i.e., every constraint token appearing in \(D_1\) appears at least as
  often in \(D_2\)).  Then \((\mathcal{C}^*/{\sim},\, \preceq)\) is a
  distributive lattice with meet \(\sqcap\) corresponding to
  coordinate-wise \(\min\) and join \(\sqcup\) to coordinate-wise \(\max\).
\end{theorem}

\begin{proof}
  The lattice \((\N_0^k, \le)\) with component-wise \(\min\) and \(\max\) is
  a distributive lattice by a standard argument (distributivity follows from
  \(\min(a, \max(b,c)) = \max(\min(a,b), \min(a,c))\) for integers).
  The PETC map is a lattice isomorphism onto its image by
  \cref{prop:petc-inj}.
\end{proof}

% ============================================================
\section{Arithmetic Control Engine (ACE) and Spectral Drift}
\label{sec:ace}
% ============================================================

\subsection{ACE Operator}
\label{ssec:ace-op}

Let \(n \in \N\) and let \(\mathcal{S}_n\) denote the space of
\(n \times n\) real symmetric matrices.

\begin{definition}[ACE Rank-Update Operator]
  \label{def:ace-op}
  For a parameter matrix \(W \in \R^{n \times n}\), an ACE
  \emph{low-rank update} of rank \(r \le n\) is
  \[
    \ACE_r(W; U, \Sigma, V)
    \;=\; W + U \Sigma V^\top,
  \]
  where \(U \in \R^{n \times r}\), \(V \in \R^{n \times r}\),
  \(\Sigma = \diag(\sigma_1, \ldots, \sigma_r)\) with
  \(\sigma_1 \ge \cdots \ge \sigma_r \ge 0\).
\end{definition}

\subsection{Spectral Drift Constraint}
\label{ssec:spectral-drift}

\begin{definition}[Spectral Drift]
  \label{def:drift}
  For a time-indexed sequence of parameter matrices
  \(\{W^{(t)}\}_{t \ge 0}\), the \emph{spectral drift} at step \(t\) is
  \[
    \delta^{(t)} \;=\; \norm{\sigma^{(t+1)} - \sigma^{(t)}}_2,
  \]
  where \(\sigma^{(t)} = (\sigma_1^{(t)}, \ldots, \sigma_n^{(t)})\) is the
  vector of singular values of \(W^{(t)}\) in non-increasing order.
\end{definition}

\begin{theorem}[Spectral Drift Bound under ACE Update]
  \label{thm:drift-bound}
  Let \(W^{(t+1)} = \ACE_r(W^{(t)}; U, \Sigma_u, V)\).  Then
  \[
    \delta^{(t)} \;\le\; \norm{U \Sigma_u V^\top}_2
    \;=\; \sigma_1(U\Sigma_u V^\top)
    \;\le\; \norm{U}_2\,\sigma_1\,\norm{V}_2.
  \]
  In particular, if \(U\) and \(V\) have orthonormal columns
  (\(\norm{U}_2 = \norm{V}_2 = 1\)), then
  \(\delta^{(t)} \le \sigma_1\).
\end{theorem}

\begin{proof}
  By Weyl's perturbation theorem for singular values, for matrices
  \(A, E \in \R^{n \times n}\),
  \[
    \abs{\sigma_i(A + E) - \sigma_i(A)} \;\le\; \norm{E}_2,
    \qquad i = 1, \ldots, n.
  \]
  Applying with \(A = W^{(t)}\) and \(E = U\Sigma_u V^\top\):
  \[
    \delta^{(t)} = \norm{\sigma^{(t+1)} - \sigma^{(t)}}_2
    \le \sqrt{\sum_{i=1}^n \bigl(\sigma_i(W^{(t)}+E) - \sigma_i(W^{(t)})\bigr)^2}
    \le \sqrt{n}\,\norm{E}_2.
  \]
  For the tighter bound, use the component-wise Weyl inequality:
  \(\max_i\abs{\sigma_i(A+E) - \sigma_i(A)} \le \norm{E}_2\), so
  \(\delta^{(t)} \le \sqrt{n}\,\norm{E}_2\).
  The sharpest practical bound follows from the submultiplicativity of the
  spectral norm:
  \[
    \norm{U\Sigma_u V^\top}_2
    \le \norm{U}_2 \cdot \norm{\Sigma_u}_2 \cdot \norm{V^\top}_2
    = \norm{U}_2 \cdot \sigma_1 \cdot \norm{V}_2.
  \]
  When \(U,V\) have orthonormal columns, \(\norm{U}_2 = \norm{V}_2 = 1\),
  giving \(\delta^{(t)} \le \sigma_1\).
\end{proof}

\begin{corollary}[ACE Drift Budget]
  \label{cor:drift-budget}
  To enforce a governance drift budget \(\eps > 0\) (i.e.,
  \(\delta^{(t)} \le \eps\) for all \(t\)), it is sufficient to require
  \[
    \sigma_1(U \Sigma_u V^\top) \;\le\; \eps.
  \]
  With orthonormal update factors this reduces to \(\sigma_1 \le \eps\).
\end{corollary}

\subsection{Frobenius Norm Accumulation}
\label{ssec:frob}

\begin{proposition}[Cumulative Drift over \(T\) Steps]
  \label{prop:cumul-drift}
  If each ACE step satisfies \(\delta^{(t)} \le \eps\) for all
  \(t \in \{0, \ldots, T-1\}\), then
  \[
    \norm{W^{(T)} - W^{(0)}}_F \;\le\; T \cdot \eps \cdot \sqrt{n}.
  \]
\end{proposition}

\begin{proof}
  By the triangle inequality and the inequality
  \(\norm{\cdot}_F \le \sqrt{n}\,\norm{\cdot}_2\):
  \[
    \norm{W^{(T)} - W^{(0)}}_F
    \le \sum_{t=0}^{T-1} \norm{W^{(t+1)} - W^{(t)}}_F
    \le \sum_{t=0}^{T-1} \sqrt{n}\,\norm{W^{(t+1)}-W^{(t)}}_2
    \le T\sqrt{n}\,\eps.
  \]
\end{proof}

% ============================================================
\section{Audit Log Coverage and Hash-Chain Integrity}
\label{sec:audit}
% ============================================================

\subsection{Coverage Metric}
\label{ssec:coverage}

Let \(\mathcal{E}\) be the universe of audit event types and let
\(\mathcal{L} = (e_1, e_2, \ldots, e_N)\) be an audit log.

\begin{definition}[Audit Coverage]
  \label{def:coverage}
  \[
    \mathrm{cov}(\mathcal{L})
    \;=\; \frac{\abs{\{e \in \mathcal{E} : \exists\, i,\; \text{type}(e_i)=e\}}}
               {\abs{\mathcal{E}}}.
  \]
\end{definition}

\begin{proposition}[Monotonicity of Coverage]
  \label{prop:cov-mono}
  If \(\mathcal{L} \subseteq \mathcal{L}'\) (as multisets), then
  \(\mathrm{cov}(\mathcal{L}) \le \mathrm{cov}(\mathcal{L}')\).
\end{proposition}

\begin{proof}
  The set \(\{\text{type}(e_i)\}\) can only grow or stay the same when more
  events are appended; division by the fixed \(\abs{\mathcal{E}}\) preserves
  the inequality.
\end{proof}

\subsection{SHA-256 Hash Chain}
\label{ssec:hashchain}

Let \(H : \{0,1\}^* \to \{0,1\}^{256}\) denote SHA-256.

\begin{definition}[Audit Hash Chain]
  \label{def:hash-chain}
  An \emph{audit hash chain} over events
  \((e_1, \ldots, e_N)\) is the sequence \((h_0, h_1, \ldots, h_N)\)
  defined recursively by
  \[
    h_0 \;=\; \mathtt{0}^{256},
    \qquad
    h_i \;=\; H\bigl(h_{i-1} \,\|\, \mathrm{enc}(e_i)\bigr),
    \quad i = 1, \ldots, N,
  \]
  where \(\|\) denotes concatenation and \(\mathrm{enc}\) is any canonical
  deterministic serialisation.
\end{definition}

\begin{theorem}[Tamper-Evidence of Hash Chain]
  \label{thm:hash-tamper}
  Under the collision-resistance assumption for \(H\), no
  probabilistic polynomial-time adversary can produce
  \((e_1',\ldots,e_N') \ne (e_1,\ldots,e_N)\) such that the
  corresponding chain terminates at the same value \(h_N\), except with
  negligible probability in the security parameter \(\kappa\).
\end{theorem}

\begin{proof}
  Suppose for contradiction that an adversary modifies event \(e_j\) to
  \(e_j' \ne e_j\) and produces the same terminal hash \(h_N\).
  By collision resistance of \(H\), the adversary cannot find
  \(h_j' = h_j\) unless it finds a collision.  A collision on step \(j\)
  propagates a different value \(h_j' \ne h_j\) into step \(j+1\), and by
  induction no subsequent value in the chain can agree with the original chain
  without finding further collisions at every remaining step.  Therefore the
  probability of producing the same \(h_N\) is at most \(N \cdot \epsilon_H\),
  where \(\epsilon_H\) is the collision-finding advantage, which is negligible
  by assumption.
\end{proof}

% ============================================================
\section{Policy State Machine: Reachability and Invariants}
\label{sec:psm}
% ============================================================

\subsection{Formal Definition}
\label{ssec:psm-def}

\begin{definition}[Policy State Machine]
  \label{def:psm}
  A \emph{Policy State Machine} (PSM) is a tuple
  \[
    \mathcal{M} \;=\; (Q, \Sigma_{\mathrm{ev}}, \delta, q_0, F, I),
  \]
  where:
  \begin{itemize}[nosep]
    \item \(Q\) is a finite set of states
          (\(\mathtt{draft},\, \mathtt{active},\, \mathtt{deprecated},\,
          \mathtt{archived},\, \mathtt{killed}\)),
    \item \(\Sigma_{\mathrm{ev}}\) is a finite event alphabet (governance
          actions),
    \item \(\delta : Q \times \Sigma_{\mathrm{ev}} \rightharpoonup Q\) is the
          partial transition function,
    \item \(q_0 = \mathtt{draft}\) is the initial state,
    \item \(F = \{\mathtt{archived}, \mathtt{killed}\}\) is the set of
          absorbing terminal states,
    \item \(I\) is a set of invariant predicates over \(Q\).
  \end{itemize}
\end{definition}

\begin{theorem}[Terminal Absorption]
  \label{thm:absorption}
  For every \(q \in F\) and every event \(a \in \Sigma_{\mathrm{ev}}\),
  \(\delta(q, a)\) is undefined; equivalently, once the PSM enters a terminal
  state it remains there.
\end{theorem}

\begin{proof}
  By construction of \(\delta\) in \texttt{engine/policy\_state\_machine.py},
  all outgoing transitions from \(F\) are omitted.  The proof is an inductive
  check on the state graph: label each node with the set of reachable states
  from it; nodes in \(F\) reach only themselves.
\end{proof}

\subsection{Kill-Switch Audit Emission Invariant}
\label{ssec:kill-switch}

\begin{proposition}[Kill-Switch Completeness]
  \label{prop:kill-switch}
  Every path in \(\mathcal{M}\) that ends in \(\mathtt{killed}\) must
  pass through exactly one transition labelled \(\mathtt{KILL}\), and that
  transition unconditionally emits an \(\mathtt{AuditEvent}\) of type
  \(\mathtt{POLICY\_KILLED}\) before the state change is committed.
\end{proposition}

\begin{proof}
  \(\mathtt{killed}\) has in-degree exactly 1 in the state graph (only the
  transition \(\mathtt{active} \xrightarrow{\mathtt{KILL}} \mathtt{killed}\)
  exists).  The implementation in \texttt{engine.py}
  wraps the state mutation in a transaction that first appends to the
  \texttt{AuditLog}; the append is synchronous and throws on failure,
  preventing state commitment if the audit write fails.
\end{proof}

% ============================================================
\section{MKT Braid Composition and Operator Norm}
\label{sec:mkt}
% ============================================================

\subsection{Braid Groups and MKT Functor}
\label{ssec:braid}

Let \(B_n\) denote the Artin braid group on \(n\) strands, generated by
\(\{\sigma_1, \ldots, \sigma_{n-1}\}\) subject to the braid relations:
\begin{align}
  \sigma_i \sigma_j &= \sigma_j \sigma_i,
  && \abs{i-j} \ge 2,\label{eq:br1}\\
  \sigma_i \sigma_{i+1} \sigma_i &= \sigma_{i+1} \sigma_i \sigma_{i+1},
  && 1 \le i \le n-2.\label{eq:br2}
\end{align}

\begin{definition}[MKT Prime Labelling]
  \label{def:mkt-label}
  The \emph{MKT prime labelling} assigns to generator \(\sigma_i\) the prime
  \(p_i\), and to its inverse \(\sigma_i^{-1}\) the formal inverse
  \(p_i^{-1}\) in \(\Q_{>0}\).  For a braid word
  \(\beta = \sigma_{i_1}^{\eps_1} \cdots \sigma_{i_k}^{\eps_k}\)
  (\(\eps_j \in \{\pm 1\}\)), define the
  \emph{MKT weight} of \(\beta\) as
  \[
    w(\beta) \;=\; \prod_{j=1}^k p_{i_j}^{\eps_j} \;\in\; \Q_{>0}.
  \]
\end{definition}

\begin{proposition}[Exponent Vector of Braid Word]
  \label{prop:braid-exponent}
  \(w(\beta) = \prod_{i=1}^{n-1} p_i^{e_i(\beta)}\), where
  \(e_i(\beta) = \sum_{j:\, i_j = i} \eps_j\) is the total signed count of
  \(\sigma_i\) appearances.  The map \(\beta \mapsto \bm{e}(\beta)\) is a
  group homomorphism from \(B_n\) to \(\Z^{n-1}\).
\end{proposition}

\begin{proof}
  Linearity of the exponent map over multiplication in \(\Q_{>0}\), combined
  with the braid relations \eqref{eq:br1}--\eqref{eq:br2} which preserve
  exponent totals, verifies that both sides transform identically under the
  relations.
\end{proof}

\subsection{Jones-MKT Operator Norm}
\label{ssec:jones-norm}

Let \(\rho : B_n \to \mathrm{GL}(V)\) be a finite-dimensional unitary
representation (e.g., the Burau representation over \(\C\)).

\begin{theorem}[Operator Norm of Braid Generator]
  \label{thm:braid-norm}
  For the reduced Burau representation at parameter \(t \in \C\),
  \(\abs{t} = 1\), the operator norm satisfies
  \[
    \norm{\rho(\sigma_i)}_2 \;=\; 1, \qquad i = 1, \ldots, n-1.
  \]
\end{theorem}

\begin{proof}
  At \(\abs{t} = 1\), the Burau matrices are unitary (this is a standard
  result; cf.\ Jones [1985]).  Unitary matrices have all singular values equal
  to 1, so \(\norm{\rho(\sigma_i)}_2 = 1\).
\end{proof}

\begin{corollary}[Submultiplicative Norm Bound for Braid Words]
  \label{cor:braid-word-norm}
  For any braid word \(\beta\) of length \(\ell\),
  \[
    \norm{\rho(\beta)}_2 \;\le\; \prod_{j=1}^{\ell}
    \norm{\rho(\sigma_{i_j}^{\eps_j})}_2 \;=\; 1.
  \]
  Equality holds iff \(\rho(\beta)\) is unitary.
\end{corollary}

\begin{proof}
  Submultiplicativity of \(\norm{\cdot}_2\) gives the inequality;
  unitarity of each generator gives the product equal to 1.
\end{proof}

\subsection{MKT Composition Stability}
\label{ssec:mkt-stability}

\begin{theorem}[MKT Functor Stability]
  \label{thm:mkt-stability}
  Let \(\beta_1, \beta_2 \in B_n\) be policy braids and let
  \(d_{\MKT}(\beta_1, \beta_2) = \norm{\bm{e}(\beta_1) - \bm{e}(\beta_2)}_1\).
  Then for any two policy braids \(\beta\) and \(\beta'\):
  \[
    d_{\MKT}(\beta \cdot \gamma,\; \beta' \cdot \gamma)
    \;=\; d_{\MKT}(\beta, \beta'),
  \]
  i.e., the metric is invariant under right-multiplication.
\end{theorem}

\begin{proof}
  Since \(\bm{e}\) is a group homomorphism,
  \(\bm{e}(\beta \cdot \gamma) = \bm{e}(\beta) + \bm{e}(\gamma)\), so
  \[
    \bm{e}(\beta\gamma) - \bm{e}(\beta'\gamma)
    = (\bm{e}(\beta) + \bm{e}(\gamma)) - (\bm{e}(\beta') + \bm{e}(\gamma))
    = \bm{e}(\beta) - \bm{e}(\beta').
  \]
  The \(\ell^1\) norm of this difference is independent of \(\gamma\).
\end{proof}

% ============================================================
\section{Lambda-M Protocol Completeness}
\label{sec:lambda}
% ============================================================

\begin{definition}[\(\Lambda^p\)-Archivum Tag]
  \label{def:lambda-tag}
  For a module \(\mathcal{A}\) with PETC encoding \(\bm{e} \in \Mspace{k}\),
  the \emph{\(\Lambda^p\)-Archivum tag} is the pair
  \[
    \mathtt{tag}(\mathcal{A}) \;=\; \bigl(\,n(\bm{e}),\; h_N(\mathcal{A})\,\bigr),
  \]
  where \(n(\bm{e}) = \prod_i p_i^{e_i}\) is the PETC numeral and
  \(h_N(\mathcal{A})\) is the terminal audit hash of \(\mathcal{A}\)'s
  audit log.
\end{definition}

\begin{theorem}[Tag Uniqueness]
  \label{thm:tag-unique}
  Under SHA-256 collision resistance, no two distinct (module, audit-log)
  pairs can share the same \(\Lambda^p\)-Archivum tag.
\end{theorem}

\begin{proof}
  Two pairs \((\mathcal{A}, \mathcal{L})\) and \((\mathcal{A}', \mathcal{L}')\)
  share a tag iff both \(n(\bm{e}(\mathcal{A})) = n(\bm{e}(\mathcal{A}'))\)
  and \(h_N(\mathcal{L}) = h_N(\mathcal{L}')\).
  The first condition forces identical multisets of constraint tokens
  (by \cref{prop:petc-inj} and unique factorisation).
  The second condition forces identical audit event sequences (by
  \cref{thm:hash-tamper}).  Together they imply
  \(\mathcal{A} = \mathcal{A}'\) and \(\mathcal{L} = \mathcal{L}'\).
\end{proof}

% ============================================================
\section{PIRTM Recursive Stability}
\label{sec:pirtm}
% ============================================================

The PIRTM (Prime Interval Recursion Theory Machine) implements a
prime-indexed rewrite system.  We state the key stability theorem.

\begin{definition}[PIRTM Rewrite Step]
  \label{def:pirtm-step}
  A PIRTM configuration is a triple \((s, \bm{e}, \tau)\) where
  \(s\) is the current state, \(\bm{e} \in \Mspace{k}\) is the working
  exponent register, and \(\tau \in \N\) is a clock.  A rewrite step
  \((s, \bm{e}, \tau) \to (s', \bm{e}', \tau+1)\) is governed by
  a prime-labelled production rule \(r = (s, p_j, \Delta\bm{e}, s')\):
  \[
    \bm{e}' \;=\; \bm{e} + \Delta\bm{e},
  \]
  where \(\Delta\bm{e}\) is a signed exponent increment in \(\Z^k\).
\end{definition}

\begin{theorem}[PIRTM Norm Contraction under Bounded Increments]
  \label{thm:pirtm-contraction}
  Suppose every production rule satisfies
  \(\norm{\Delta\bm{e}}_1 \le \delta_{\max}\).  Then after \(T\) steps,
  \[
    \norm{\bm{e}^{(T)} - \bm{e}^{(0)}}_1 \;\le\; T \cdot \delta_{\max}.
  \]
  If moreover \(\delta_{\max} < \norm{\bm{e}^{(0)}}_1 / T\), the system
  cannot exhaust its initial exponent budget in \(T\) steps.
\end{theorem}

\begin{proof}
  By the triangle inequality applied to the telescoping sum:
  \[
    \norm{\bm{e}^{(T)} - \bm{e}^{(0)}}_1
    \le \sum_{t=0}^{T-1} \norm{\bm{e}^{(t+1)} - \bm{e}^{(t)}}_1
    = \sum_{t=0}^{T-1} \norm{\Delta\bm{e}^{(t)}}_1
    \le T \cdot \delta_{\max}.
  \]
  The second claim follows directly.
\end{proof}

% ============================================================
\section{Galois Certification and Field Extensions}
\label{sec:galois}
% ============================================================

\begin{definition}[Policy Galois Group]
  \label{def:galois}
  Let \(f(x) = \prod_{p \in S} (x - p) \in \Z[x]\) for a finite prime set
  \(S \subset \primes\).  The \emph{policy Galois group}
  \(\mathrm{Gal}(f/\Q)\) acts on \(S\) by permuting roots.
  For distinct primes, \(f\) factors into distinct linear factors over \(\Q\),
  giving \(\mathrm{Gal}(f/\Q) \cong \{e\}\) (trivial).
\end{definition}

\begin{proposition}[Trivial Policy Galois Group]
  \label{prop:galois-trivial}
  For any finite set of distinct rational primes \(S\),
  \(\mathrm{Gal}(f/\Q) \cong \{e\}\).
\end{proposition}

\begin{proof}
  Every prime \(p \in S\) is rational, so all roots of \(f\) lie in \(\Q\)
  itself.  The splitting field is \(\Q\), and
  \(\mathrm{Gal}(\Q/\Q) = \{e\}\).
\end{proof}

\begin{remark}
  This triviality is intentional: prime-labelled policy tokens are
  \emph{absolutely distinguishable} with no Galois symmetry; there is no
  automorphism swapping two distinct policy constraints.  This is the
  algebraic expression of the uniqueness requirement in governance semantics.
\end{remark}

% ============================================================
\section{Agent Scope Containment Theorem}
\label{sec:agent}
% ============================================================

\begin{definition}[Agent Scope Vector]
  \label{def:scope}
  An agent \(\mathcal{G}\) has a \emph{scope vector}
  \(\bm{s}_{\mathcal{G}} \in \{0,1\}^m\), where \(m = |\mathcal{C}|\),
  indicating which constraint tokens it is authorised to invoke.
\end{definition}

\begin{definition}[Scope-Bounded Execution]
  \label{def:scope-bounded}
  An execution trace \((c_1, \ldots, c_\ell)\) is
  \emph{scope-bounded for} \(\mathcal{G}\) iff
  \(\PETC(c_1,\ldots,c_\ell) \le \bm{s}_{\mathcal{G}}\)
  component-wise (where inequality is in \(\N_0^m\) and
  \(\bm{s}_\mathcal{G}\) is treated as an indicator-budget: at most once per
  authorised token unless explicitly budget-expanded).
\end{definition}

\begin{theorem}[Scope Containment]
  \label{thm:scope}
  If every agent transition is checked against the scope predicate
  \(\PETC(\text{trace}) \le \bm{s}_{\mathcal{G}}\) before commitment, then
  no execution trace can invoke an unauthorised constraint token.
\end{theorem}

\begin{proof}
  By contrapositive: if an unauthorised token \(c^* \notin \mathrm{supp}(\bm{s}_\mathcal{G})\)
  were invoked, the corresponding coordinate of \(\PETC(\text{trace})\) would
  become \(\ge 1 > 0 = (\bm{s}_\mathcal{G})_{c^*}\), violating the predicate.
  The pre-commitment check enforces this as a precondition, so such a
  transition is rejected.
\end{proof}

% ============================================================
\section{Summary of Norm Bounds}
\label{sec:summary}
% ============================================================

\renewcommand{\arraystretch}{1.3}
\begin{table}[ht]
\centering
\caption{Operator Norm Bounds Established in This Appendix}
\label{tab:bounds}
\begin{tabular}{@{}llll@{}}
\toprule
\textbf{Object} & \textbf{Norm} & \textbf{Bound} & \textbf{Reference} \\
\midrule
ACE rank-\(r\) update \(U\Sigma V^\top\)
  & \(\norm{\cdot}_2\)
  & \(\le \sigma_1 \norm{U}_2 \norm{V}_2\)
  & \cref{thm:drift-bound} \\
ACE cumulative drift (\(T\) steps)
  & \(\norm{\cdot}_F\)
  & \(\le T\eps\sqrt{n}\)
  & \cref{prop:cumul-drift} \\
Burau generator \(\rho(\sigma_i)\), \(\abs{t}=1\)
  & \(\norm{\cdot}_2\)
  & \(= 1\)
  & \cref{thm:braid-norm} \\
Braid word \(\rho(\beta)\), length \(\ell\)
  & \(\norm{\cdot}_2\)
  & \(\le 1\)
  & \cref{cor:braid-word-norm} \\
PIRTM \(\ell^1\)-drift (\(T\) steps)
  & \(\norm{\cdot}_1\)
  & \(\le T\delta_{\max}\)
  & \cref{thm:pirtm-contraction} \\
PETC distance (multisets)
  & \(d_{\PETC}\)
  & \(\norm{\cdot}_1\) metric
  & \cref{def:petc-dist} \\
\bottomrule
\end{tabular}
\end{table}

% ============================================================
\section{Code Correspondence}
\label{sec:code}
% ============================================================

The following listing shows the Python implementation of the ACE drift check
as implemented in \texttt{debug\_spectral.py}, corresponding to
\cref{thm:drift-bound} and \cref{cor:drift-budget}.

\begin{lstlisting}[language=Python, caption={ACE Spectral Drift Enforcement (\texttt{debug\_spectral.py})}, label={lst:ace-drift}]
import numpy as np
from typing import Optional

def ace_rank_update(
    W: np.ndarray,
    U: np.ndarray,
    sigma_vals: np.ndarray,
    V: np.ndarray,
    eps: float,
) -> Optional[np.ndarray]:
    """
    Apply ACE low-rank update W' = W + U diag(sigma_vals) V^T,
    enforcing spectral drift constraint:
        sigma_1(U diag(sigma_vals) V^T) <= eps.
    Returns updated W' if constraint holds, else None.

    Corresponds to Theorem 3.2 (drift bound) and Corollary 3.3
    (drift budget) in the Mathematical Appendix.
    """
    Sigma_u = np.diag(sigma_vals)
    E = U @ Sigma_u @ V.T                    # low-rank perturbation
    drift = np.linalg.norm(E, ord=2)         # spectral norm = sigma_1(E)
    if drift > eps:
        return None                          # reject: budget exceeded
    W_new = W + E
    # Verify Weyl bound on singular value shift
    sv_old = np.linalg.svd(W,     compute_uv=False)
    sv_new = np.linalg.svd(W_new, compute_uv=False)
    delta  = np.linalg.norm(sv_new - sv_old, ord=2)
    assert delta <= drift + 1e-10, (          # Weyl: delta <= ||E||_2
        f"Weyl bound violated: delta={delta:.6f}, drift={drift:.6f}"
    )
    return W_new
\end{lstlisting}

The following listing implements the PETC subsumption check from
\cref{thm:petc-subsumption} and \cref{thm:scope-containment}.

\begin{lstlisting}[language=Python, caption={PETC Scope Check (\texttt{engine/predicate\_registry.py})}, label={lst:petc-scope}]
from sympy import factorint
from typing import Dict

PRIME_INDEX: Dict[str, int] = {
    # token -> prime assignment
    "read":   2,  "write":  3, "execute": 5,
    "delete": 7,  "admin": 11, "audit":  13,
}

def petc_encode(tokens: list[str]) -> Dict[int, int]:
    """Return exponent vector as {prime: exponent} dict."""
    result: Dict[int, int] = {}
    for tok in tokens:
        p = PRIME_INDEX[tok]
        result[p] = result.get(p, 0) + 1
    return result

def scope_check(trace_tokens: list[str], scope_tokens: list[str]) -> bool:
    """
    Enforce Theorem 8.3: no trace token may exceed scope budget.
    Returns True iff PETC(trace) <= PETC(scope) component-wise.
    """
    trace_ev = petc_encode(trace_tokens)
    scope_ev = petc_encode(scope_tokens)
    for prime, count in trace_ev.items():
        if scope_ev.get(prime, 0) < count:
            return False   # unauthorised token invoked
    return True
\end{lstlisting}

% ============================================================
\section{Open Problems and Future Work}
\label{sec:open}
% ============================================================

\begin{enumerate}[leftmargin=2em, label=\textbf{O\arabic*.}]
  \item \textbf{Tighter Weyl Bound.}
    \cref{thm:drift-bound} uses the standard Weyl perturbation bound,
    which is tight in the worst case but may be loose for structured
    low-rank updates.  Derive tighter bounds exploiting the rank-\(r\)
    structure directly via the Eckart--Young theorem.

  \item \textbf{Knot Invariant Integration.}
    Extend the MKT functor (\cref{sec:mkt}) to full knot invariants
    (HOMFLY-PT, Khovanov homology) and characterise which governance
    policies correspond to slice knots.

  \item \textbf{Post-Quantum Hash Chain.}
    Replace SHA-256 in \cref{def:hash-chain} with a lattice-based hash
    function (e.g., SWIFFT) and re-establish \cref{thm:hash-tamper} under
    the SIVP hardness assumption.

  \item \textbf{Dirichlet L-Function Policy Encoding.}
    Encode policy documents as characters
    \(\chi : (\Z/m\Z)^\times \to \C^\times\) and represent the PETC
    product as \(L(s, \chi)\); explore whether analytic continuation
    provides additional expressivity for temporal policy semantics.

  \item \textbf{PIRTM Termination Decidability.}
    Characterise the class of PIRTM rule sets for which termination is
    decidable, in analogy with Minsky machine halting results.
\end{enumerate}

% ============================================================
%  END OF APPENDIX
% ============================================================

\vspace{1.5em}
\noindent\rule{\linewidth}{0.4pt}\\[0.3em]
{\small
\noindent\textbf{Defensive Publication Notice.}
This document constitutes a prior art defensive publication establishing
priority date 25 April 2026 for all mathematical constructions,
theorems, proofs, operator norm bounds, and algorithm designs described
herein.  All rights reserved under the Prime Materia License v1.0
and Apache License 2.0, copyright \copyright\ 2026 Multiplicity Foundation /
Ryan Van Gelder.  Unrestricted academic and non-commercial use is permitted
with attribution.
\nocite{*}
\bibliographystyle{unsrt}  
\bibliography{references}
\end{document}